% !TEX program = pdflatex
\documentclass[11pt,a4paper]{article}

\usepackage[utf8]{inputenc}
\usepackage[T1]{fontenc}
\usepackage[ngerman]{babel}
\usepackage{lmodern}
\usepackage[margin=2.4cm]{geometry}
\usepackage{amsmath,amssymb,amsthm}
\usepackage{booktabs}
\usepackage{enumitem}
\usepackage{microtype}
\usepackage{hyperref}
\usepackage{xcolor}
\usepackage{array}

\hypersetup{
  colorlinks=true,
  linkcolor=black,
  citecolor=black,
  urlcolor=blue!50!black
}

\newtheorem{theorem}{Satz}
\newtheorem{proposition}[theorem]{Proposition}
\newtheorem{lemma}[theorem]{Lemma}
\theoremstyle{definition}
\newtheorem{definition}[theorem]{Definition}
\newtheorem{hypothesis}[theorem]{Hypothese}
\theoremstyle{remark}
\newtheorem{remark}[theorem]{Bemerkung}

\newcommand{\Lean}{\textsc{Lean}\,4}
\newcommand{\CA}{\ensuremath{[\mathrm{C}\!\to\!\mathrm{A}]}}
\newcommand{\tagA}{\ensuremath{[\mathrm{A}]}}
\newcommand{\tagB}{\ensuremath{[\mathrm{B}]}}
\newcommand{\tagC}{\ensuremath{[\mathrm{C}]}}
% Kurznamen in Formeln; Lean-Identifikatoren im Fließtext.
\newcommand{\FeedIn}{\mathrm{FeedIn}}
\newcommand{\Residual}{\mathrm{Residual}}
\newcommand{\FiniteExit}{\mathrm{FiniteExit}}
\newcommand{\Forever}{\mathrm{ForeverExp}}
\newcommand{\Umap}{U}
\newcommand{\N}{\mathbb{N}}
\newcommand{\AP}{\mathrm{AP}}

\title{%
  Von der Collatz-Struktur zur Residualdynamik\\[0.35em]
  \large Das Core6-Formalisierungsprogramm:\\
  historischer Kontext, algebraische Reduktion\\
  und die verbleibende dynamische Barriere
}
\author{%
  Thomas Hofbauer\\
  \small Bamberg / Akademie Olympia\\
  \small\texttt{Kepler--Hurwitz}-Formalisierungsprogramm
}
\date{27.\ Juli 2026}

\begin{document}
\maketitle

\begin{abstract}
Die Collatz-Vermutung bleibt eines der hartnäckigsten offenen Probleme der
diskreten Mathematik. Seit ihrer Formulierung durch Lothar Collatz in den
1930er Jahren widersteht sie Ansätzen von elementaren Kongruenzargumenten
über probabilistische Dynamik bis hin zu großflächiger Rechnung.

Dieser Artikel stellt ein Formalisierungsprogramm auf der Basis der
Core6-Zylinderarchitektur vor. Ziel ist \emph{nicht} der Anspruch eines
Beweises der Collatz-Vermutung, sondern die präzise Isolation jener Teile
des Problems, die algebraisch zugänglich werden, und jenes Teils, der
genuin dynamisch bleibt.

Das Programm trennt drei grundsätzlich verschiedene Schichten:
\begin{enumerate}[leftmargin=*,itemsep=0.2em]
  \item die arithmetische Struktur von Bewertungszylindern,
  \item endliche Residuenraum-Besetzung und Dichteaussagen,
  \item globale Orbit-Erreichbarkeit.
\end{enumerate}
Im Core6-Rahmen werden feste Bewertungsfasern klassifiziert, die
kontrahierende Residuenbesetzung durch ein dyadisches Zertifikat
quantifiziert und endliche Core6-interne First-Hit-Pfade als
arithmetische Progressionen beschrieben.

Eine stärkere universelle Pfadüberdeckungsbehauptung wird durch einen
expliziten Zeugen formal widerlegt. Die verbleibende Collatz-Schwierigkeit
reduziert sich damit auf ein Residual-Erreichbarkeitsproblem.

Die zentrale Schlussfolgerung ist daher keine Lösung von Collatz, sondern
ein Lokalisierungssatz:
\[
  \text{globale Erreichbarkeit}
  \;\Longleftrightarrow\;
  \text{Erreichbarkeit der Residualmenge}.
\]
Die arithmetische Struktur ist klassifiziert; die dynamische Barriere bleibt.
Epistemischer Status: \texttt{ClaimsFreeze} ist falsch; die Collatz-Vermutung
bleibt unbewiesen.
\end{abstract}

\tableofcontents
\vspace{1em}

%==============================================================================
\section{Historischer Hintergrund: von Collatz zur modernen Formaldynamik}
%==============================================================================

Die Iteration, die heute als Collatz-Abbildung bekannt ist, wurde von
Lothar Collatz im zwanzigsten Jahrhundert eingeführt. Das Problem fragt,
ob die wiederholte Anwendung von
\[
  T(n)
  \;=\;
  \begin{cases}
    n/2, & n\equiv 0\pmod{2},\\
    3n+1, & n\equiv 1\pmod{2}
  \end{cases}
\]
schließlich den Zyklus
\[
  4\to 2\to 1
\]
erreicht. Obwohl die Regel elementar ist, verbinden die resultierende
Dynamik lokale arithmetische Einfachheit mit globaler Unvorhersehbarkeit.

Frühe Untersuchungen von Stanisław Ulam und Shizuo Kakutani erkannten das
Problem als diskretes dynamisches System. Everett und andere entwickelten
probabilistische Sichtweisen; daraus entstand das moderne Verständnis, dass
mittleres Verhalten wesentlich leichter zugänglich sein kann als eine
vollständige Orbitklassifikation.

Riho Terras begründete zentrale Dichte- und Stoppzeitresultate. Seine Arbeit
führte eine methodisch unverzichtbare Unterscheidung ein:
\[
  \text{fast-alles Verhalten}
  \;\neq\;
  \text{alles Verhalten}.
\]
Diese Unterscheidung bleibt zentral. Probabilistische Abstiegsmodelle können
typisches Verhalten stark nahelegen und zugleich Ausnahmebahnen vollständig
offenlassen. Wenn im Projekt informell von ``Mr.~Taro'' die Rede ist, ist
Terras gemeint: die Stimme, die Heuristik von Zertifikat trennt.

Jeffrey Lagarias systematisierte weite Teile der bekannten Theorie und
betonte Kongruenzstrukturen, rechnerische Verifikation sowie die präzise
Trennung zwischen bewiesenen Aussagen und Heuristiken.

John Conways Resultate zu verallgemeinerten Collatz-Systemen zeigten, dass
verwandte Probleme sogar unentscheidbares Verhalten kodieren können; das
verstärkt die Notwendigkeit mathematischer Vorsicht.

Spätere Arbeiten von Wirsching, Korec und vielen rechnerischen Schulen
erweiterten das strukturelle und experimentelle Verständnis der Iteration.
Terence Taos Ergebnis von 2019, dass fast alle Bahnen fast beschränkte Werte
erreichen, ist ein bedeutender Meilenstein, bleibt aber ein Dichtesatz und
kein Beweis universeller Termination~\cite{Tao2019}.

Das vorliegende Programm folgt dieser historischen Lehre:
\begin{quote}
\emph{Eine Reduktion ist nur dann wertvoll, wenn die verbleibende logische
Lücke explizit benannt wird.}
\end{quote}

%==============================================================================
\section{Formale Methodik}
%==============================================================================

Das Projekt nutzt \Lean{} als Beweisassistenten. Jede mathematische Aussage
trägt einen epistemischen Status:

\begin{center}
\begin{tabular}{@{}llp{7.2cm}@{}}
\toprule
Tag & Bedeutung & Aufstiegsregel \\
\midrule
\tagA & formal bewiesener Satz & Repository of record nach Review \\
\tagB & Census / Numerik & hebt allein niemals zu \tagA{} auf \\
\CA & konstruktives Lean-Paket, noch nicht \tagA & Merge + separater Promotion-Commit \\
\tagC & offene Hypothese / Interface & kann zu \CA{} werden \\
\bottomrule
\end{tabular}
\end{center}

Diese Unterscheidung verhindert einen typischen Fehlmodus der
rechnerischen Mathematik: umfangreiche Verifikation mit Beweis zu
verwechseln. Ziel ist daher nicht bloß die Erzeugung von Daten, sondern die
Umwandlung struktureller Beobachtungen in kerngeprüfte Mathematik.

Global bleibt \texttt{ClaimsFreeze} \textbf{falsch}: es wird kein Collatz-Beweis
beansprucht.

%==============================================================================
\section{Die Core6-Zylinderarchitektur}
%==============================================================================

\subsection{Die ungerade Syracuse-Abbildung}

Eine oft klarere Formulierung arbeitet mit der ungeraden Syracuse-Abbildung
\[
  \Umap(n)
  \;=\;
  \frac{3n+1}{2^{v_2(3n+1)}},
\]
dem ersten ungeraden Iterat nach einem ungeraden Kick. Erreichbarkeit des
trivialen Zyklus ist dann äquivalent dazu, dass jeder ungerade Start unter
Iteraten von~$\Umap$ schließlich in ein kontrahierendes Regime eintritt.

\subsection{Faserworte und der Core6-Präfix}

Ein \emph{Faserwort} ist eine endliche Folge positiver $2$-adischer
Bewertungen, die die sukzessiven Halbierungen nach ungeraden Syracuse-Kicks
aufzeichnet. Der ausgezeichnete Präfix
\[
  \mathtt{core6}
  \;=\;
  [1,1,1,1,2,2]
\]
ist ein Bewertungswort der Länge~$6$, das von einer kohärenten Familie von
Starts geteilt wird. Anhängen eines positiven Schwanzexponenten~$e$ liefert
die Länge-$7$-Faser
\[
  \mathtt{fiberE}(e)
  \;=\;
  \mathtt{core6}\mathbin{+\kern-0.15em+}[e].
\]

\subsection{Kanonische Basen und Zylinder}

Für jeden Exponenten $e\ge 1$ gibt es einen eindeutigen
\emph{kanonischen Repräsentanten} $b_e$ in einem vollständigen
Restsystem modulo $2^{e+9}$ und eine unendliche arithmetische Progression---den
\emph{kanonischen Zylinder}
\[
  C_e
  \;=\;
  \bigl\{\, b_e + k\cdot 2^{e+9}
  \;\big|\;
  k\in\N \,\bigr\}.
\]
Eine konzeptionell unverzichtbare Typentrennung lautet:
\[
  b_e
  \;\neq\;
  [b_e]
  \;\neq\;
  C_e.
\]
Der Repräsentant ist ein endliches Objekt; die Restklasse ist ein
Äquivalenzobjekt; der Zylinder ist eine unendliche Teilmenge der natürlichen
Zahlen. Diese Trennung verhindert den Typenkollaps zwischen
arithmetischen Repräsentanten und unendlichen Fasern.

\subsection{Expandierende und kontrahierende Masse}

Dynamisch splittet man
\begin{align*}
  \mathtt{expandingMass}
  &\;=\; C_1\cup C_2\cup C_3,\\
  \mathtt{contractingMass}
  &\;=\; \textstyle\bigcup_{e\ge 4} C_e.
\end{align*}
Heuristisch tendieren expandierende Kanäle unter einem vollen Block zum
Wachstum, kontrahierende zum Schrumpfen. Das klassische Erreichbarkeitsziel
lautet in dieser Sprache:

\begin{hypothesis}[ReachabilityFeedInGoal \tagC]
Jedes $n\in\mathtt{expandingMass}$ gestattet ein $t\ge 1$ mit
$\Umap^{\circ t}(n)\in\mathtt{contractingMass}$.
\end{hypothesis}

Diese Hypothese bleibt offen. Was sich ändert, ist ihre
\emph{logische Lage}.

%==============================================================================
\section{Statische Klassifikation: PR\,\#15 und PR\,\#16}
%==============================================================================

Die erste Phase etabliert die algebraische Struktur.

\subsection{Kanonisches Samen-Lifting (PR\,\#15)}

Für $e\ge 4$ gehören die Core6-Zylinder zum kontrahierenden Regime.
Die Faserindex-Abbildung
\[
  \Phi_e(k)
  \;=\;
  b_e + k\, M_e
\]
ist injektiv und liefert eine Äquivalenz $\N\simeq C_e$. Jeder
Zylinderpunkt ist eine unendliche arithmetische Progression echter
Core6-Schwanz-Realisierungen. Algebraisch ist der ungerade affine
Koeffizient entlang $\mathtt{fiberE}(e)$ unabhängig von~$e$ (die Konstante
$2347$ in der Formalisierung).

\subsection{Statische Zylinderpartition (PR\,\#16)}

Paarweise Disjunktheit der Zylinder~$C_e$, die expandierende/kontrahierende
Phasenpartition des Core6-Zylinders und eine \emph{relative} dyadische
Dichteaussage:
\[
  \frac{\lvert R_m^{\mathrm{contr}}\rvert}
       {\lvert R_m^{\mathrm{Core6}}\rvert}
  \;=\;
  \frac18 - \frac1{2^m}.
\]
Daher
\[
  \lim_{m\to\infty}
  \frac{\lvert R_m^{\mathrm{contr}}\rvert}
       {\lvert R_m^{\mathrm{Core6}}\rvert}
  \;=\;
  \frac18.
\]
Dieses Resultat ist rein statisch. Es impliziert \textbf{nicht}:
Orbitkonvergenz, universelle Erreichbarkeit oder Collatz-Termination.
Die Unterscheidung zwischen Residuenbesetzung und dynamischer Anziehung ist
fundamental.

\begin{remark}[Expliziter Nicht-Claim]
Die relative dyadische Dichte ist \textbf{keine} natürliche Dichte auf~$\N$
und kein Collatz-Satz. Status der gelieferten Pakete: derzeit \CA{}
(CI-grüne Kandidaten), noch nicht repository-\tagA{} nach Merge und
separatem Promotion-Commit.
\end{remark}

%==============================================================================
\section{Dynamische Phase: PR\,\#17 und PR\,\#18}
%==============================================================================

Die nächste Phase untersucht, ob expandierende Zylinder in die
kontrahierende Core6-Masse einspeisen.

\subsection{Ein-Block-Feed-In und D2b-Freeze (PR\,\#17)}

Nach einem vollen $\mathtt{fiberE}(e_0)$-Block bilden jene Starts in~$C_{e_0}$,
die sofort in kontrahierende Masse landen, die Menge
$\mathtt{oneBlockFeedInSet}(e_0)$. Diese Menge ist eine disjunkte Vereinigung
arithmetischer Progressionen, indiziert durch Zielkanäle $f\ge 4$, via einer
Cancel-by-$2$-Kongruenzbrücke und einer eindeutigen Residuenklasse~$\kappa$.

Ein positiver Zeuge ist $1{,}246{,}239\in\mathtt{oneBlockFeedInSet}(1)$;
der universelle Ein-Block-Anspruch ist formal falsch via $31\mapsto 137$.

\subsection{Kanalpfade und First-Hit-Packaging (PR\,\#18 / D3)}

Die anfängliche Hypothese war, dass jeder expandierende Core6-Zylinder einen
endlichen Core6-internen First-Hit-Pfad gestatten könnte. Das wurde über
Kanalpfade formalisiert:
\[
  p
  \;=\;
  [e_0,e_1,\ldots,e_r].
\]
Für jeden \emph{festen} Pfad bilden die Realisierer eine arithmetische
Progression:
\[
  \bigl\{\, n \;\big|\; \mathtt{RealizesChannelPath}(p,n) \,\bigr\}
  \;=\;
  \AP_p.
\]
Der algebraische Mechanismus ist iteriertes dyadisches Kongruenzlifting,
nicht ein klassisches CRT als Blackbox. Die Pfadstruktur ist damit
vollständig klassifiziert.

Ein \emph{kontrahierender First-Hit-Pfad} startet in einem expandierenden
Kanal, bleibt in jedem inneren Kanal expandierend und endet beim ersten
kontrahierenden Kanal. Ihre Vereinigung ist in Lean
\texttt{core6PathFeedInSet}$(e_0)$; in Formeln schreiben wir kurz
\[
  \FeedIn(e_0)
  \;=\;
  \bigcup_{\substack{p\text{ kontrahierend}\\\text{First-Hit von }e_0}}
  \AP_p.
\]
Das ist der Core6-\emph{interne} erfolgreiche Kanal.

\begin{theorem}[First-Hit-Brücke \CA]
Ist $n\in\FeedIn(e_0)$, so existieren $r\ge 1$ und $t=7r$ mit
$\Umap^{\circ t}(n)\in\mathtt{contractingMass}$.
\end{theorem}
Der Core6-interne Anteil braucht daher \textbf{keine weitere offene Dynamik}.

%==============================================================================
\section{Das Scheitern der universellen Core6-Pfadüberdeckung}
%==============================================================================

Die stärkere Aussage
\[
  \forall\, n\in C_1\cup C_2\cup C_3\;
  \exists\, p:\ \mathtt{RealizesChannelPath}(p,n)
\]
(in der präzisen Form $\mathtt{Core6PathFeedInGoal}$: jeder Start in~$C_{e_0}$
realisiert einen Core6-internen kontrahierenden First-Hit-Pfad) ist falsch.

\begin{theorem}[$\neg\,\mathtt{Core6PathFeedInGoal}$ \CA]
Es ist nicht wahr, dass jeder Start in $C_1\cup C_2\cup C_3$ einen
Core6-internen kontrahierenden First-Hit-Pfad realisiert.
\end{theorem}

Ein formaler Gegenzeuge ist
\[
  31\in C_1,
  \qquad
  31\notin\FeedIn(1),
\]
denn nach einem Block gilt $31\mapsto 137\notin\mathtt{core6Cylinder}$.
Jeder Core6-interne First-Hit-Pfad würde das erste Blockbild zurück in
Core6 zwingen.

\begin{remark}
Dieses Scheitern zerstört die algebraische Klassifikation nicht. Es
identifiziert die fehlende dynamische Komponente. Die Strategie
``jeder expandierende Start hat einen rein Core6-internen First-Hit-Pfad''
ist nicht bloß unvollendet---sie ist \textbf{falsch}.
\end{remark}

%==============================================================================
\section{Residualreduktion}
%==============================================================================

\begin{definition}[Residualmenge]
In Lean: \texttt{core6PathResidualSet}$(e_0)$.
\[
  \Residual(e_0)
  \;=\;
  C_{e_0}\setminus\FeedIn(e_0).
\]
\end{definition}

\begin{theorem}[Zylinderdichotomie \CA]
\[
  C_{e_0}
  \;=\;
  \FeedIn(e_0)
  \;\dot\cup\;
  \Residual(e_0),
\]
und $\mathtt{oneBlockOffCore6Set}(e_0)\subseteq\Residual(e_0)$.
\end{theorem}

Für $e_0=1$ ist das Residual nichtleer ($31$) und der Feed-In-Anteil
nichtleer ($1{,}246{,}239$). Der Feed-In-Anteil ist über arithmetische
Progressionskanäle verstanden; das Residual enthält genau jene Startwerte,
die vom endlichen Core6-internen Pfadmechanismus nicht erfasst werden.

\begin{hypothesis}[ResidualFeedInGoal \tagC]
Jedes Residual-Start erreicht unter ungerader Syracuse-Iteration schließlich
kontrahierende Masse.
\end{hypothesis}

\begin{theorem}[Residualreduktion \CA]
\[
  \boxed{
  \mathtt{ResidualFeedInGoal}
  \;\Longleftrightarrow\;
  \mathtt{ReachabilityFeedInGoal}
  }
\]
\end{theorem}

\begin{remark}
Die Rückrichtung ist unmittelbar: Residualstarts liegen in expandierender
Masse. Die Vorwärtsrichtung braucht die Dichotomie \emph{und} die
First-Hit-Brücke. Die bloße Mengengleichung
$C=F\,\dot\cup\,R$ verschiebt für sich allein keine Dynamik von~$F$ in die
kontrahierende Masse.
\end{remark}

In einem Satz:
\begin{quote}
\emph{Alles, was innerhalb Core6-interner First-Hit-Pfade abgeschlossen werden
kann, ist abgeschlossen; alles, was bleibt, ist genau das Residual---und das
Residual ist so schwer wie die gesamte Expanding-Mass-Erreichbarkeitsfrage.}
\end{quote}

Mathematischer Freeze-Kopf von D3.3:
\texttt{24bd5c83718da11c9ec2bfa8962de331be117369}.

%==============================================================================
\section{Die D3.4-Forschungsgrenze: Residualfeinstruktur}
%==============================================================================

Das Residual ist nicht homogen. Die nächste strukturelle Zerlegung lautet:
\[
  \Residual(e_0)
  \;=\;
  \FiniteExit(e_0)
  \;\dot\cup\;
  \Forever(e_0).
\]

\begin{definition}[FiniteBlockExit]
Lean-Name: \texttt{finiteBlockExitSet}.
$\FiniteExit(e_0)$ umfasst Starts mit expandierendem Core6-Aufenthalt, deren
nächster voller Block den Core6-Zylinder verlässt.
Tiefe~$0$ ist \texttt{oneBlockOffCore6Set}$(e_0)$; Zeuge: $31\in\FiniteExit(1)$.
\end{definition}

\begin{definition}[ForeverExpandingBoundary]
Lean-Name: \texttt{foreverExpandingBoundarySet}.
$\Forever(e_0)$ umfasst Starts, deren jeder expandierende Aufenthalt beim
nächsten Block wieder in expandierender Masse landet.
\end{definition}

\begin{theorem}[Residualfeinstruktur \CA\ (D3.4)]
Für expandierende Startkanäle $e_0\in\{1,2,3\}$ gilt
\[
  \Residual(e_0)
  \;=\;
  \FiniteExit(e_0)
  \;\dot\cup\;
  \Forever(e_0),
\]
mit $\FiniteExit(e_0)\cap\Forever(e_0)=\emptyset$,
$\mathtt{oneBlockOffCore6Set}(e_0)\subseteq\FiniteExit(e_0)
\subseteq\Residual(e_0)$ und
$\Forever(e_0)\subseteq\Residual(e_0)$.
\end{theorem}

Diese Unterscheidung trennt zwei grundsätzlich verschiedene Fragen:
\begin{enumerate}[leftmargin=*]
  \item endliche arithmetische Fluchtmechanismen
        ($\FiniteExit$, inkl.\ möglicher späterer Re-Entry),
  \item unendliche dynamische Obstruktion
        ($\Forever$).
\end{enumerate}
Beide dynamischen Ziele bleiben \tagC:
$\mathtt{OffCore6ReentryGoal}$ und die Ausschließung bzw.\ Kontrolle von
$\Forever$ sind nicht Bestandteil der algebraischen Dichotomie.
Insbesondere ist $\mathtt{OffCore6ReentryGoal}$ \textbf{kein} Ersatz für
$\mathtt{ResidualFeedInGoal}$.

%==============================================================================
\section{Governance: wie Überanspruch vermieden wird}
%==============================================================================

\begin{enumerate}[leftmargin=*]
  \item PR\,\#18 steht unter strengem mathematischem Freeze bei
  \texttt{24bd5c8}. Neue Kernel-Mathematik gehört nicht auf diesen Pull Request.
  \item D3.4/D4 leben auf einem Follow-up-Branch
  (\texttt{cursor/core6-dynamic-d34-residual-\ldots}).
  \item Status-Tags bleiben \CA{} bis Review, Merge und einem
  \textbf{separaten} Promotion-Commit zu repository-\tagA.
  \item Natürliche Dichte, globales Collatz und ``Collapse''-Rhetorik sind
  explizite Nicht-Ziele des vorliegenden Pakets.
\end{enumerate}

Diese Disziplin ist die Schuld an Conways Unentscheidbarkeitswarnung und an
Lagarias' Survey-Kultur: den Satz benennen, den man hat---und den Satz
benennen, den man nicht hat.

%==============================================================================
\section{Abschlussstatus}
%==============================================================================

\begin{center}
\small
\begin{tabular}{@{}p{4.8cm}p{2.6cm}p{4.6cm}@{}}
\toprule
Komponente & Status & Kommentar \\
\midrule
Core6-Algebra / Faserklassifikation
  & bewiesen (\CA/\tagA-Pfad) & PR\,\#15--\#16 \\
Relative dyadische Dichte $1/8$
  & bewiesen (\CA) & kein Natürlich-Dichte-Claim \\
Feste Kanalpfade $=\AP_p$
  & bewiesen (\CA) & D3.2 \\
Universelle Core6-Pfadüberdeckung
  & widerlegt & $\neg$\,Core6PathFeedInGoal \\
Feed-In / Residual-Zerlegung
  & bewiesen (\CA) & D3.3c \\
Residual $\Leftrightarrow$ Reachability
  & bewiesen (\CA) & D3.3d + First-Hit-Brücke \\
Residual-Feinstruktur
  & bewiesen (\CA)
  & $\FiniteExit\,\dot\cup\,\Forever$ \\
Residual-Erreichbarkeit
  & offen \tagC & D4 \\
Collatz-Vermutung
  & offen & \texttt{ClaimsFreeze} false \\
\bottomrule
\end{tabular}
\end{center}

Die korrekte Schlussfolgerung lautet daher:
\[
  \boxed{\text{Die arithmetische Struktur ist klassifiziert.}}
\]
\[
  \boxed{\text{Die verbleibende Schwierigkeit ist genau Residualdynamik.}}
\]
Kein Beweis von Collatz wird beansprucht. Der Beitrag ist eine formal
verifizierte Reduktion dahin, \emph{wo} ein solcher Beweis noch operieren muss.

\vspace{1.2em}
\begin{flushright}
\textit{Bamberg, 27.\ Juli 2026}\\
Thomas Hofbauer\\
\texttt{Kepler--Hurwitz / Pre119--Core6}
\end{flushright}

%==============================================================================
\begin{thebibliography}{9}

\bibitem{Collatz}
L.~Collatz.
\newblock Zum $3x+1$-Problem (historische Zuschreibung; zirkulierende Notizen
und Korrespondenz, Mitte des 20.~Jahrhunderts).
\newblock Siehe Surveys in~\cite{Lagarias2010}.

\bibitem{Terras1976}
R.~Terras.
\newblock A stopping time problem on the positive integers.
\newblock \emph{Acta Arith.} \textbf{30} (1976), 241--252.

\bibitem{Lagarias2010}
J.~C.~Lagarias (Hrsg.).
\newblock \emph{The Ultimate Challenge: The $3x+1$ Problem}.
\newblock American Mathematical Society, 2010.

\bibitem{Conway1972}
J.~H.~Conway.
\newblock Unpredictable iterations.
\newblock In \emph{Proc.\ 1972 Number Theory Conf.}, Univ.\ of Colorado,
Boulder, 1972, 49--52.

\bibitem{Tao2019}
T.~Tao.
\newblock Almost all orbits of the Collatz map attain almost bounded values.
\newblock arXiv:1909.03562, 2019.

\bibitem{Wirsching1998}
G.~J.~Wirsching.
\newblock \emph{The Dynamical System Generated by the $3n+1$ Function}.
\newblock Lecture Notes in Mathematics \textbf{1681}, Springer, 1998.

\bibitem{KHPre119}
Akademie Olympia / Kepler--Hurwitz.
\newblock Pre119-Core6-Formalisierungsstack (PRs \#15--\#18 und D3.4-Follow-up),
Math-Freeze-Köpfe
\texttt{3740183} (statisch),
\texttt{18e8747} (D2b),
\texttt{24bd5c8} (D3.3);
\texttt{ClaimsFreeze} falsch.

\end{thebibliography}

\end{document}
